% ============================================================
% CHƯƠNG 4: THỰC NGHIỆM VÀ KẾT QUẢ
% ============================================================

\chapter{Thực nghiệm và kết quả}
\label{chap:experiments}

\section{Môi trường thực nghiệm}
\label{sec:experiment_env}

Toàn bộ thực nghiệm được tiến hành trên một máy tính cá nhân với cấu hình phần cứng và phần mềm được trình bày trong Bảng~\ref{tab:hardware_software}.

\begin{table}[H]
\centering
\caption{Cấu hình môi trường thực nghiệm}
\label{tab:hardware_software}
\small
\begin{tabular}{|l|l|}
\hline
\textbf{Thành phần} & \textbf{Chi tiết} \\
\hline
\multicolumn{2}{|c|}{\textit{Phần cứng}} \\
\hline
CPU               & Intel Core (multi-core) \\
RAM               & 16 GB \\
GPU               & Không sử dụng (huấn luyện trên CPU) \\
\hline
\multicolumn{2}{|c|}{\textit{Phần mềm}} \\
\hline
Hệ điều hành     & Windows 10/11 \\
Python            & 3.12.4 \\
TensorFlow        & 2.16.1 \\
NumPy             & 1.26.4 \\
scikit-learn      & 1.5.x \\
SHAP              & 0.46.0 \\
Keras Tuner       & 1.4.x \\
pandas            & 2.2.x \\
matplotlib        & 3.9.x \\
\hline
\end{tabular}
\end{table}

Lưu ý kỹ thuật: TensorFlow 2.16.1 yêu cầu NumPy $< 2.0$ để đảm bảo tương thích. SHAP 0.46.0 cũng yêu cầu NumPy $< 2.0$. Toàn bộ quá trình huấn luyện được thực hiện trên CPU, do đó thời gian huấn luyện dài hơn so với GPU nhưng kết quả mô hình không bị ảnh hưởng.

% ============================================================

\section{Kết quả nhận dạng hoạt động (HAR)}
\label{sec:har_results}

Module HAR (Mục~\ref{sec:har_module}) được huấn luyện trên tập WISDM v1.1 với kiến trúc Stacked Bi-LSTM (2 lớp $\times$ 30 units/chiều). Kết quả phân loại trên tập test đạt accuracy tổng thể \textbf{96.19\%}.

\begin{table}[H]
\centering
\caption{Kết quả phân loại HAR trên tập test WISDM v1.1}
\label{tab:har_test_results}
\small
\begin{tabular}{|l|c|c|c|c|}
\hline
\textbf{Hoạt động} & \textbf{Precision} & \textbf{Recall} & \textbf{F1-Score} & \textbf{Samples} \\
\hline
Downstairs & 0.86 & 0.82 & 0.84 & 471 \\
Jogging    & 0.97 & 0.99 & 0.98 & 1{,}933 \\
Sitting    & 1.00 & 0.91 & 0.95 & 306 \\
Standing   & 0.90 & 1.00 & 0.95 & 246 \\
Upstairs   & 0.85 & 0.84 & 0.84 & 632 \\
Walking    & 0.97 & 0.98 & 0.97 & 2{,}081 \\
\hline
\textbf{Weighted Avg} & \textbf{0.96} & \textbf{0.96} & \textbf{0.96} & \textbf{5{,}669} \\
\hline
\end{tabular}
\end{table}

\textbf{Nhận xét:}
\begin{itemize}
    \item \textbf{Jogging} dễ nhận dạng nhất (F1 = 0.98) — tín hiệu gia tốc kế có biên độ cao và nhịp điệu rõ ràng.
    \item \textbf{Walking} cũng đạt F1 cao (0.97) — mẫu bước đi đặc trưng ở cả 3 trục.
    \item \textbf{Downstairs/Upstairs} khó phân biệt hơn (F1 = 0.84) — tín hiệu leo cầu thang có nhiều điểm tương đồng với đi bộ.
    \item \textbf{Sitting/Standing} đạt F1 = 0.95 mặc dù cả hai là hoạt động tĩnh — mô hình nắm bắt được sự khác biệt tinh tế trong dao động gia tốc kế.
\end{itemize}

Mô hình HAR đạt accuracy 96.19\% đủ tốt để cung cấp nhãn hoạt động chính xác cho module sinh dữ liệu. Accuracy cao ở Jogging (99\% recall) và Walking (98\% recall) — hai hoạt động phổ biến nhất — đảm bảo chất lượng dữ liệu đa phương thức.

% ============================================================

\section{Kết quả mô hình Baseline (Stacked Bi-LSTM)}
\label{sec:baseline_results}

Mô hình Stacked Bi-LSTM Baseline (128$\to$64 units, dropout=0.3, lr=0.001) được huấn luyện trên tập 13 đặc trưng. Kết quả trên tập test:

\begin{table}[H]
\centering
\caption{Kết quả mô hình Baseline trên tập test}
\label{tab:baseline_test_metrics}
\begin{tabular}{|l|c|}
\hline
\textbf{Metric} & \textbf{Giá trị} \\
\hline
MAE  & 0.6855 \\
RMSE & 0.8723 \\
R²   & 0.9245 \\
MSE  & 0.7609 \\
\hline
\end{tabular}
\end{table}

R² = 0.9245 cho thấy mô hình giải thích được 92.45\% phương sai của mức stress. MAE = 0.6855 nghĩa là trung bình mỗi dự đoán lệch khoảng 0.69 điểm trên thang 1--9 — mức sai số chấp nhận được cho ứng dụng theo dõi xu hướng stress.

Tuy nhiên, mô hình baseline còn dư địa cải thiện: dropout = 0.3 có thể quá cao (underfitting nhẹ), và learning rate = 0.001 có thể khiến quá trình tối ưu hội tụ chậm hoặc mắc kẹt tại local minimum. Đây là động lực cho bước tối ưu siêu tham số tiếp theo.

% ============================================================

\section{Kết quả tối ưu siêu tham số}
\label{sec:hp_tuning_results}

Bayesian Optimization được thực hiện với 20 trials trên không gian tìm kiếm đã định nghĩa (Bảng~\ref{tab:search_space}). Bảng~\ref{tab:best_hp} trình bày siêu tham số tốt nhất so với baseline.

\begin{table}[H]
\centering
\caption{Siêu tham số tốt nhất từ Bayesian Optimization}
\label{tab:best_hp}
\begin{tabular}{|l|c|c|c|}
\hline
\textbf{Tham số} & \textbf{Baseline} & \textbf{Tuned} & \textbf{Thay đổi} \\
\hline
\texttt{lstm\_units\_1}  & 128   & \textbf{64}    & $\downarrow$ 50\% \\
\texttt{lstm\_units\_2}  & 64    & \textbf{64}    & -- \\
\texttt{dropout\_rate}   & 0.3   & \textbf{0.1}   & $\downarrow$ 67\% \\
\texttt{dense\_units}    & 64    & \textbf{128}   & $\uparrow$ 100\% \\
\texttt{learning\_rate}  & 0.001 & \textbf{0.01}  & $\uparrow$ 10$\times$ \\
\hline
Tổng params              & 320.129 & \textbf{163.585} & $\downarrow$ 48.9\% \\
\hline
\end{tabular}
\end{table}

\textbf{Kết quả sau tối ưu:}

\begin{table}[H]
\centering
\caption{Cải thiện hiệu suất sau Hyperparameter Tuning}
\label{tab:tuning_improvement}
\begin{tabular}{|l|c|c|c|c|}
\hline
\textbf{Metric} & \textbf{Baseline} & \textbf{Tuned} & \textbf{Thay đổi} & \textbf{\%} \\
\hline
MAE  & 0.6855 & \textbf{0.5292} & $-$0.1563 & $\mathbf{-22.8\%}$ \\
RMSE & 0.8723 & \textbf{0.7483} & $-$0.1240 & $\mathbf{-14.2\%}$ \\
R²   & 0.9245 & \textbf{0.9444} & $+$0.0199 & $\mathbf{+2.2\%}$ \\
\hline
\end{tabular}
\end{table}

\textbf{Phân tích ý nghĩa:}

\begin{itemize}
    \item \textbf{Dropout 0.3 $\to$ 0.1}: Baseline bị underfitting nhẹ — giảm dropout cho phép mô hình khai thác nhiều neurons hơn khi huấn luyện, tăng khả năng nắm bắt mẫu phức tạp.
    \item \textbf{Learning rate 0.001 $\to$ 0.01}: Tốc độ học cao hơn giúp mô hình hội tụ nhanh hơn và thoát local minima. Kết hợp với ReduceLROnPlateau (giảm 50\% khi plateau), learning rate sẽ tự động giảm dần về giá trị nhỏ hơn ở cuối quá trình huấn luyện.
    \item \textbf{Tổng params giảm 48.9\%} (320K $\to$ 164K): Mô hình tuned nhỏ hơn đáng kể nhưng hiệu suất tốt hơn — cho thấy baseline bị \textit{over-parameterized} (quá nhiều tham số so với độ phức tạp dữ liệu).
\end{itemize}

% ============================================================

\section{Phân tích tầm quan trọng đặc trưng}
\label{sec:feature_importance_results}

Để hiểu đóng góp của từng đặc trưng vào dự đoán stress, khóa luận áp dụng 4 phương pháp phân tích bổ sung nhau:

\begin{enumerate}
    \item \textbf{Permutation Importance}: Đo mức giảm hiệu suất (MAE) khi xáo trộn ngẫu nhiên từng đặc trưng \cite{fisher2019permutation}.
    \item \textbf{SHAP (SHapley Additive exPlanations)}: Sử dụng KernelExplainer tính giá trị Shapley trung bình cho từng đặc trưng \cite{lundberg2017shap}.
    \item \textbf{Correlation Analysis}: Tính hệ số tương quan Pearson giữa từng đặc trưng và Stress\_Level.
    \item \textbf{RF Surrogate}: Huấn luyện Random Forest trên dữ liệu đã chuẩn hóa, dùng Gini importance \cite{breiman2001random}.
\end{enumerate}

\begin{table}[H]
\centering
\caption{Top-5 đặc trưng theo 4 phương pháp phân tích}
\label{tab:feature_importance_top5}
\small
\begin{tabular}{|c|l|l|l|l|}
\hline
\textbf{Rank} & \textbf{Permutation} & \textbf{SHAP} & \textbf{Correlation} & \textbf{RF Surrogate} \\
\hline
1 & Heart\_Rate       & Heart\_Rate       & Mood\_Score   & Location \\
2 & Mood\_Score        & Mood\_Score        & Phone\_Event  & Mood\_Score \\
3 & Screen\_Usage      & Screen\_Usage      & Heart\_Rate   & Phone\_Event \\
4 & Energy\_Level      & Day\_of\_Week      & Location      & Heart\_Rate \\
5 & Day\_of\_Week      & Hour              & Screen\_Usage & Energy\_Level \\
\hline
\end{tabular}
\end{table}

Hình~\ref{fig:feature_importance} minh họa mức độ quan trọng của 13 đặc trưng theo 4 phương pháp.

\begin{figure}[H]
    \centering
    \includegraphics[width=\linewidth]{../results/feature_importance_13features/feature_importance_13features.png}
    \caption{Tầm quan trọng đặc trưng theo 4 phương pháp phân tích}
    \label{fig:feature_importance}
\end{figure}

Bảng~\ref{tab:overall_ranking} tổng hợp thứ hạng tổng thể (trung bình hạng từ 4 phương pháp).

\begin{table}[H]
\centering
\caption{Xếp hạng tổng thể 13 đặc trưng (trung bình hạng từ 4 phương pháp)}
\label{tab:overall_ranking}
\small
\begin{tabular}{|c|l|c|c|c|c|c|}
\hline
\textbf{Rank} & \textbf{Đặc trưng} & \textbf{Perm.} & \textbf{SHAP} & \textbf{Corr.} & \textbf{RF} & \textbf{TB hạng} \\
\hline
1  & Mood\_Score              & 2  & 2  & 1  & 2  & 1.75 \\
2  & Heart\_Rate              & 1  & 1  & 3  & 4  & 2.25 \\
3  & Screen\_Usage\_Current   & 3  & 3  & 5  & 6  & 4.25 \\
4  & Phone\_Event\_Frequency  & 6  & 6  & 2  & 3  & 4.25 \\
5  & Location                 & 8  & 8  & 4  & 1  & 5.25 \\
6  & Day\_of\_Week            & 5  & 4  & 6  & 9  & 6.00 \\
7  & Energy\_Level            & 4  & 7  & 8  & 5  & 6.00 \\
8  & Hour                     & 7  & 5  & 7  & 8  & 6.75 \\
9  & Sleep\_Duration          & 10 & 9  & 13 & 7  & 9.75 \\
10 & Activity                 & 9  & 11 & 11 & 12 & 10.75 \\
11 & Accelerometer\_Y         & 11 & 10 & 12 & 10 & 10.75 \\
12 & Accelerometer\_Z         & 12 & 13 & 10 & 11 & 11.50 \\
13 & Accelerometer\_X         & 13 & 12 & 9  & 13 & 11.75 \\
\hline
\end{tabular}
\end{table}

Hình~\ref{fig:ranking_comparison} so sánh thứ hạng đặc trưng giữa các phương pháp.

\begin{figure}[H]
    \centering
    \includegraphics[width=\linewidth]{../results/feature_importance_13features/ranking_comparison_13features.png}
    \caption{So sánh thứ hạng đặc trưng giữa 4 phương pháp}
    \label{fig:ranking_comparison}
\end{figure}

\textbf{Nhận xét:}

\begin{enumerate}
    \item \textbf{Mood\_Score và Heart\_Rate nhất quán đứng top-2}: Cả 4 phương pháp đều xếp hai đặc trưng này ở vị trí đầu (hạng trung bình 1.75 và 2.25). Kết quả phù hợp với nghiên cứu: Hovsepian et al. \cite{hovsepian2015cstress} xác nhận nhịp tim là chỉ số stress sinh lý quan trọng nhất, còn tâm trạng (mood) phản ánh trực tiếp trạng thái tâm lý.
    
    \item \textbf{Screen\_Usage\_Current (hạng 3)} và \textbf{Phone\_Event\_Frequency (hạng 4)}: Hai đặc trưng hành vi liên quan đến sử dụng điện thoại chiếm vị trí quan trọng, gợi ý rằng mẫu sử dụng smartphone có liên hệ mạnh với mức stress — phù hợp với xu hướng nghiên cứu digital phenotyping.
    
    \item \textbf{Location quan trọng ở RF nhưng thấp ở Permutation/SHAP}: Location chiếm 67.58\% Gini importance trong RF (hạng 1) nhưng chỉ hạng 8 ở Permutation. Điều này do RF dùng Decision Tree — dễ khai thác biến categorical với nhiều lớp (6 locations), trong khi LSTM (Permutation/SHAP) nắm bắt ảnh hưởng phức tạp hơn thông qua interaction với các đặc trưng khác.
    
    \item \textbf{Accelerometer features có importance thấp nhất} (hạng 11--13): Mặc dù cảm biến gia tốc là đầu vào chính của module HAR, thông tin chuyển động thô ít trực tiếp ảnh hưởng đến stress. Thông tin hoạt động đã được ``tóm tắt'' qua nhãn Activity, nên Accelerometer trở thành ``đặc trưng dự phòng''.
\end{enumerate}

% ============================================================

\section{Phân tích lỗi}
\label{sec:error_analysis}

\subsection{Phân bổ lỗi tổng thể}
\label{subsec:error_distribution}

Bảng~\ref{tab:error_distribution} trình bày phân bổ lỗi tổng thể của mô hình baseline và tuned.

\begin{table}[H]
\centering
\caption{Phân bổ lỗi dự đoán: Baseline vs Tuned}
\label{tab:error_distribution}
\small
\begin{tabular}{|l|c|c|}
\hline
\textbf{Chỉ số} & \textbf{Baseline} & \textbf{Tuned} \\
\hline
Mean Error (bias)     & $+$0.295  & $-$0.150 \\
Std Error             & 0.821     & 0.733 \\
MAE                   & 0.685     & 0.529 \\
Median AE             & 0.547     & 0.340 \\
P90                   & 1.391     & 1.101 \\
P95                   & 1.685     & 1.555 \\
P99                   & 2.200     & 2.418 \\
Max Error             & 5.352     & 5.047 \\
\% $<$ 0.5 điểm       & 47.7\%   & \textbf{60.3\%} \\
\% $<$ 1.0 điểm       & 73.2\%   & \textbf{82.7\%} \\
\% $<$ 2.0 điểm       & 97.3\%   & \textbf{97.7\%} \\
\hline
\end{tabular}
\end{table}

Mô hình tuned có 82.7\% dự đoán sai dưới 1.0 điểm (so với 73.2\% baseline). Bias giảm từ $+$0.295 (baseline thiên dự đoán cao) xuống $-$0.150 (tuned thiên dự đoán thấp nhẹ). Đặc biệt, median AE giảm từ 0.547 xuống 0.340 — phần lớn dự đoán chỉ lệch khoảng 1/3 điểm.

Hình~\ref{fig:error_comparison} minh họa so sánh lỗi trực quan giữa baseline và tuned.

\begin{figure}[H]
    \centering
    \includegraphics[width=\linewidth]{../results/error_analysis_13features_tuned/error_comparison_baseline_vs_tuned.png}
    \caption{So sánh phân bổ lỗi giữa mô hình Baseline và Tuned}
    \label{fig:error_comparison}
\end{figure}

\subsection{Phân tích lỗi theo mức stress}
\label{subsec:error_by_stress}

Bảng~\ref{tab:error_by_stress} phân tích MAE theo 4 nhóm mức stress.

\begin{table}[H]
\centering
\caption{MAE theo mức stress: Baseline vs Tuned}
\label{tab:error_by_stress}
\small
\begin{tabular}{|l|c|c|c|c|c|}
\hline
\textbf{Mức stress} & \textbf{Samples} & \textbf{Base MAE} & \textbf{Tuned MAE} & \textbf{Thay đổi} & \textbf{\%} \\
\hline
Low (1--3)       & 4{,}221 & 0.319 & 0.341 & $+$0.022 & $+$6.9\% \\
Medium (4--5)    & 1{,}236 & 0.874 & 0.751 & $-$0.124 & $\mathbf{-14.1\%}$ \\
High (6--7)      & 573     & 0.907 & 0.903 & $-$0.004 & $-$0.4\% \\
Very High (8--9) & 2{,}078 & 1.257 & 0.677 & $-$0.580 & $\mathbf{-46.1\%}$ \\
\hline
\end{tabular}
\end{table}

\textbf{Nhận xét:}
\begin{itemize}
    \item \textbf{Very High stress (8--9) cải thiện mạnh nhất} ($-$46.1\%): Baseline dự đoán nhóm này rất kém (MAE = 1.257, bias = $+$1.257 — luôn dự đoán thấp hơn thực tế). Tuned giảm MAE xuống 0.677, gần với các nhóm khác.
    \item \textbf{Low stress (1--3) tăng nhẹ} ($+$6.9\%): MAE tăng không đáng kể (0.319 $\to$ 0.341). Nguyên nhân: tuned ``đánh đổi'' một chút chính xác ở nhóm dễ để cải thiện nhóm khó (Very High).
    \item \textbf{Medium stress (4--5) vẫn là nhóm khó nhất}: MAE = 0.751 — mức stress trung bình nằm ở vùng chuyển tiếp, nơi nhiều đặc trưng có giá trị gần nhau, khiến mô hình khó phân biệt.
\end{itemize}

\subsection{Phân tích lỗi theo hoạt động}
\label{subsec:error_by_activity}

\begin{table}[H]
\centering
\caption{MAE theo hoạt động: Baseline vs Tuned}
\label{tab:error_by_activity}
\small
\begin{tabular}{|l|c|c|c|c|}
\hline
\textbf{Hoạt động} & \textbf{Samples} & \textbf{Base MAE} & \textbf{Tuned MAE} & \textbf{Thay đổi (\%)} \\
\hline
Upstairs    & 467   & 0.950 & 0.727 & $-$23.5\% \\
Downstairs  & 416   & 0.913 & 0.792 & $-$13.3\% \\
Sitting     & 2{,}337 & 0.731 & 0.502 & $\mathbf{-31.2\%}$ \\
Jogging     & 814   & 0.648 & 0.514 & $-$20.6\% \\
Walking     & 2{,}258 & 0.639 & 0.525 & $-$17.8\% \\
Standing    & 1{,}816 & 0.582 & 0.465 & $-$20.2\% \\
\hline
\end{tabular}
\end{table}

\textbf{Nhận xét:}
\begin{itemize}
    \item \textbf{Sitting cải thiện mạnh nhất} ($-$31.2\%): Đây là hoạt động phổ biến (2.337 mẫu) với biến động stress lớn (ngồi ở office vs ngồi ở nhà có stress rất khác). Tuning giúp mô hình nắm bắt tốt hơn ngữ cảnh sitting.
    \item \textbf{Standing có MAE thấp nhất} (0.465): Hoạt động đứng thường xảy ra ở ngữ cảnh ổn định, ít biến động → dễ dự đoán.
    \item \textbf{Upstairs/Downstairs có MAE cao nhất}: Ít mẫu (467, 416) và tín hiệu sensor biến động mạnh.
    \item Tất cả hoạt động đều được cải thiện bởi tuning.
\end{itemize}

\subsection{Phân tích lỗi theo thời gian}
\label{subsec:error_by_time}

\begin{table}[H]
\centering
\caption{MAE theo khoảng thời gian trong ngày: Baseline vs Tuned}
\label{tab:error_by_time}
\small
\begin{tabular}{|l|c|c|c|c|}
\hline
\textbf{Khoảng thời gian} & \textbf{Samples} & \textbf{Base MAE} & \textbf{Tuned MAE} & \textbf{Thay đổi (\%)} \\
\hline
Night (0--5)       & --      & --    & --    & -- \\
Morning (6--11)    & 2{,}109 & 0.847 & 0.650 & $-$23.3\% \\
Afternoon (12--17) & 3{,}024 & 0.877 & 0.689 & $-$21.4\% \\
Evening (18--23)   & 2{,}975 & 0.376 & 0.281 & $\mathbf{-25.4\%}$ \\
\hline
\end{tabular}
\end{table}

\textbf{Nhận xét:}
\begin{itemize}
    \item Không có mẫu ban đêm (0--5h) trong tập test vì dữ liệu mô phỏng lịch trình ngủ đêm (0--6h).
    \item \textbf{Evening có MAE thấp nhất} (0.281): Buổi tối thường ở nhà, hoạt động ổn định, stress giảm dần — mẫu dễ học.
    \item \textbf{Morning và Afternoon có MAE cao hơn}: Giai đoạn hoạt động chính trong ngày (làm việc, di chuyển, tập thể dục) có biến động stress lớn hơn.
    \item Tuning cải thiện đồng đều ở mọi khoảng thời gian ($-$21\% đến $-$25\%).
\end{itemize}

% ============================================================

\section{So sánh các mô hình}
\label{sec:model_comparison_results}

5 kiến trúc mô hình được huấn luyện trên cùng pipeline dữ liệu, cùng callbacks, cùng batch size (32), tối đa 80 epochs. Bảng~\ref{tab:model_comparison} trình bày kết quả tổng hợp.

\begin{table}[H]
\centering
\caption{So sánh 5 mô hình trên tập test}
\label{tab:model_comparison}
\small
\begin{tabular}{|c|l|c|c|c|c|c|c|}
\hline
\textbf{\#} & \textbf{Mô hình} & \textbf{MAE} & \textbf{RMSE} & \textbf{R²} & \textbf{Params} & \textbf{Time (s)} & \textbf{Epochs} \\
\hline
1 & MLP (Dense)              & 0.931 & 1.297 & 0.833 & 241K & 60    & 19 \\
2 & Simple LSTM              & 0.521 & 0.760 & 0.943 & 83K  & 452   & 17 \\
3 & Bi-LSTM (Baseline)       & 0.716 & 0.970 & 0.907 & 320K & 1{,}148 & 16 \\
4 & Stacked Bi-GRU           & 0.755 & 0.910 & 0.918 & 244K & 2{,}855 & 46 \\
5 & \textbf{Bi-LSTM (Tuned)} & \textbf{0.441} & \textbf{0.670} & \textbf{0.956} & \textbf{164K} & 825 & 20 \\
\hline
\end{tabular}
\end{table}

Hình~\ref{fig:model_comparison} trực quan hóa kết quả so sánh các metrics chính.

\begin{figure}[H]
    \centering
    \includegraphics[width=\linewidth]{../results/model_comparison/model_comparison.png}
    \caption{So sánh hiệu suất 5 mô hình}
    \label{fig:model_comparison}
\end{figure}

Hình~\ref{fig:model_radar} cung cấp góc nhìn radar chart đa chiều.

\begin{figure}[H]
    \centering
    \includegraphics[width=0.8\linewidth]{../results/model_comparison/model_comparison_radar.png}
    \caption{Radar chart so sánh đa chiều 5 mô hình}
    \label{fig:model_radar}
\end{figure}

\textbf{Phân tích chi tiết:}

\begin{enumerate}
    \item \textbf{MLP có hiệu suất thấp nhất} (R² = 0.833, MAE = 0.931):
    
    MLP flatten toàn bộ chuỗi thành vector phẳng, mất hoàn toàn thông tin thứ tự thời gian. Kết quả R² = 0.833 so với Simple LSTM R² = 0.943 cho thấy \textit{temporal dependency đóng góp $\sim$11\% R²} — chứng minh rằng thứ tự thời gian của dữ liệu stress là thông tin quan trọng, không chỉ là giá trị tĩnh của các đặc trưng.
    
    \item \textbf{Simple LSTM (R² = 0.943) bất ngờ tốt hơn Bi-LSTM Baseline (R² = 0.907):}
    
    Kết quả này có vẻ trái trực giác — mô hình 1 lớp, 1 chiều, ít tham số hơn (83K vs 320K) lại tốt hơn mô hình phức tạp hơn. Giải thích: Bi-LSTM Baseline có quá nhiều tham số (over-parameterized) với dropout cao (0.3) — sự kết hợp này gây \textit{underfitting}: mô hình không đủ ``tự do'' để nắm bắt mẫu. Simple LSTM nhỏ gọn hơn, dropout thấp hơn (0.3 chỉ ở 1 lớp), nên thực tế ``fit'' dữ liệu tốt hơn.
    
    Phát hiện này là \textit{bằng chứng thực nghiệm mạnh} cho tầm quan trọng của hyperparameter tuning: kiến trúc tốt (Bi-LSTM) với siêu tham số xấu có thể thua kiến trúc đơn giản hơn (Simple LSTM) với siêu tham số phù hợp.
    
    \item \textbf{Stacked Bi-GRU (R² = 0.918) $\approx$ Bi-LSTM Baseline (R² = 0.907):}
    
    GRU và LSTM cho kết quả gần nhau khi cùng cấu hình — phù hợp với nhận xét của Chung et al. \cite{chung2014empirical} rằng hai kiến trúc có hiệu suất tương đương trên nhiều bài toán. Tuy nhiên, Bi-GRU huấn luyện 46 epochs (lâu nhất, 2.855s) do hội tụ chậm hơn.
    
    \item \textbf{Stacked Bi-LSTM Tuned tối ưu nhất} (R² = 0.956, MAE = 0.441):
    
    Mô hình tuned vượt trội toàn diện:
    \begin{itemize}
        \item MAE thấp nhất (0.441) — dự đoán stress chỉ lệch trung bình 0.44 điểm trên thang 1--9.
        \item R² cao nhất (0.956) — giải thích 95.6\% phương sai.
        \item Params vừa phải (164K) — ít hơn Baseline (320K) và MLP (241K).
        \item Thời gian huấn luyện hợp lý (825s) — nhanh hơn Bi-GRU (2.855s) và Baseline (1.148s).
    \end{itemize}
    
    \item \textbf{Bài học từ so sánh}: Hyperparameter tuning quan trọng hơn việc thay đổi kiến trúc. So sánh mô hình 3 vs 5 (cùng Bi-LSTM, chỉ khác HP):
    \begin{itemize}
        \item MAE: 0.716 $\to$ 0.441 ($\mathbf{-38.4\%}$)
        \item R²: 0.907 $\to$ 0.956 ($+5.4\%$)
        \item Params: 320K $\to$ 164K ($-48.9\%$)
    \end{itemize}
    Chỉ bằng cách \textit{tìm siêu tham số phù hợp}, cùng kiến trúc Bi-LSTM đã cải thiện 38.4\% MAE và giảm gần nửa số tham số.
\end{enumerate}

% ============================================================

\section{Thảo luận theo câu hỏi nghiên cứu}
\label{sec:rq_discussion}

Phần này tổng hợp kết quả thực nghiệm theo các câu hỏi nghiên cứu đã nêu ở Chương~\ref{chap:introduction}.

\textbf{RQ1 -- Context-aware có giúp diễn giải tín hiệu sinh lý tốt hơn không?}

Kết quả phân tích đặc trưng cho thấy các biến ngữ cảnh-hành vi (Mood\_Score, Screen\_Usage\_Current, Phone\_Event\_Frequency, Location) cùng với Heart\_Rate đồng thời nằm trong nhóm đóng góp cao. Điều này cho thấy mô hình không dựa vào sinh lý đơn biến mà học theo tổ hợp ngữ cảnh. Tuy nhiên, để khẳng định quan hệ nhân quả chặt chẽ hơn, nghiên cứu tiếp theo nên bổ sung thí nghiệm ablation trực tiếp (bật/tắt Context-Stress Modifier).

\textbf{RQ2 -- Mô hình chuỗi thời gian có vượt trội mô hình phi chuỗi không?}

Kết quả ủng hộ rõ ràng: MLP (R² = 0.833) thấp hơn đáng kể so với các mô hình hồi quy theo chuỗi như Simple LSTM (R² = 0.943) và Bi-LSTM Tuned (R² = 0.956). Do đó, giả thuyết về vai trò của phụ thuộc thời gian được xác nhận.

\textbf{RQ3 -- Bayesian Optimization có tạo khác biệt đáng kể so với baseline không?}

Kết quả xác nhận mạnh: từ Bi-LSTM Baseline sang Bi-LSTM Tuned, MAE giảm 38.4\%, R² tăng 5.4\%, đồng thời số tham số giảm 48.9\%. Điều này cho thấy đóng góp của tối ưu siêu tham số là bản chất, không chỉ là tinh chỉnh nhỏ.

\textbf{RQ4 -- Nhóm đặc trưng nào đóng góp chính cho dự đoán stress?}

Bốn phương pháp giải thích độc lập đều chỉ ra sự nhất quán cao ở nhóm sinh lý-hành vi, đặc biệt là Mood\_Score và Heart\_Rate. Kết quả này củng cố giả thuyết rằng dự đoán stress hiệu quả cần tích hợp đa nguồn tín hiệu thay vì chỉ dựa trên dữ liệu cảm biến chuyển động.

% ============================================================

\section{So sánh với nghiên cứu liên quan}
\label{sec:related_comparison}

Bảng~\ref{tab:related_work_comparison} so sánh kết quả của khóa luận với các công trình đã công bố trong lĩnh vực dự đoán/phát hiện stress.

\begin{table}[H]
\centering
\caption{So sánh với các nghiên cứu liên quan}
\label{tab:related_work_comparison}
\small
\begin{tabular}{|l|l|l|l|l|}
\hline
\textbf{Nghiên cứu} & \textbf{Phương pháp} & \textbf{Bài toán} & \textbf{Dữ liệu} & \textbf{Kết quả} \\
\hline
Hovsepian et al. \cite{hovsepian2015cstress} & SVM + ECG/RIP & Binary & cStress & F1 = 0.72 \\
Schmidt et al. \cite{schmidt2018wesad} & Random Forest & 3-class & WESAD & Acc = 0.80 \\
Can et al. \cite{can2019stress} & CNN + wearable & Binary & Custom & Acc = 0.92 \\
Rashid et al. \cite{rashid2024stress} & Bi-LSTM & Multi-class & Custom & F1 = 0.85 \\
\hline
\textbf{Nghiên cứu này} & \textbf{Stacked Bi-LSTM} & \textbf{Regression (1--9)} & \textbf{Custom} & $\mathbf{R^2 = 0.956}$ \\
\hline
\end{tabular}
\end{table}

\textbf{Nhận xét:}

\begin{enumerate}
    \item \textbf{Bài toán khác biệt}: Phần lớn nghiên cứu hiện tại giải quyết bài toán phân loại (binary hoặc multi-class), trong khi khóa luận này giải quyết bài toán \textit{hồi quy} (regression) — dự đoán mức stress liên tục trên thang 1--9. Bài toán regression khó hơn classification (dự đoán giá trị chính xác thay vì chỉ phân lớp) và cung cấp thông tin chi tiết hơn về mức độ stress.
    
    \item \textbf{So sánh gián tiếp}: Do khác bài toán và dữ liệu, không thể so sánh trực tiếp metrics (F1 vs R²). Tuy nhiên, R² = 0.956 cho thấy mô hình nắm bắt được 95.6\% biến động stress — đây là kết quả mạnh cho bài toán regression.
    
    \item \textbf{Đóng góp riêng}: Khóa luận này khác biệt ở:
    \begin{itemize}
        \item Sử dụng dữ liệu đa phương thức (multi-modal): cảm biến gia tốc kế + sinh lý + hành vi + ngữ cảnh.
        \item Tích hợp HAR vào pipeline stress prediction (context-aware).
        \item So sánh có hệ thống 5 kiến trúc, phân tích tầm quan trọng đặc trưng bằng 4 phương pháp.
        \item Bayesian Optimization tối ưu siêu tham số — cải thiện 38.4\% MAE.
    \end{itemize}
    
    \item \textbf{Hạn chế}: Dữ liệu của khóa luận là dữ liệu mô phỏng (simulated), không phải thu thập thực tế. Do đó, kết quả R² cao có thể phần nào phản ánh tính nhất quán của mô hình sinh dữ liệu. Để khẳng định mạnh hơn, cần kiểm chứng trên dữ liệu thực tế.
\end{enumerate}

% ============================================================

\section{Đe dọa đến tính hiệu lực của nghiên cứu}
\label{sec:validity_threats}

Để giữ trọng tâm của Chương 4 vào báo cáo và đối chiếu kết quả thực nghiệm, phần phân tích chi tiết các đe dọa đến tính hiệu lực (construct, internal, external, conclusion validity) được chuyển sang Mục~\ref{sec:limitations} của Chương~\ref{chap:conclusion} dưới dạng các hạn chế nghiên cứu.

Trong chương này, kết quả được diễn giải theo phạm vi dữ liệu và thiết kế thí nghiệm đã nêu; các giới hạn suy luận và hướng kiểm chứng bổ sung được thảo luận tập trung ở Chương 5.

% ============================================================

\section{Tổng kết chương}
\label{sec:chapter4_summary}

Chương này trình bày kết quả thực nghiệm toàn diện của hệ thống dự đoán stress:

\begin{enumerate}
    \item \textbf{HAR Module}: Accuracy 96.19\% trên WISDM — cung cấp nhãn hoạt động chính xác.
    \item \textbf{Baseline Bi-LSTM}: R² = 0.9245, MAE = 0.6855 — hiệu suất tốt ngay từ đầu.
    \item \textbf{Bayesian Optimization}: Cải thiện $-$22.8\% MAE chỉ bằng tối ưu siêu tham số, giảm 48.9\% tham số.
    \item \textbf{Feature Importance}: Heart\_Rate và Mood\_Score nhất quán là 2 đặc trưng quan trọng nhất ở cả 4 phương pháp.
    \item \textbf{Error Analysis}: Tuned cải thiện mạnh nhất ở Very High stress ($-$46.1\%) và Sitting ($-$31.2\%).
    \item \textbf{Model Comparison}: Stacked Bi-LSTM Tuned đạt R² = 0.9555, MAE = 0.4414 — tốt nhất trong 5 kiến trúc.
    \item \textbf{Bài học chính}: Hyperparameter tuning quan trọng hơn thay đổi kiến trúc — cùng Bi-LSTM, tuning cải thiện 38.4\% MAE.
\end{enumerate}
